\documentclass[10pt,a4paper]{amsart}
\usepackage[margin=19mm]{geometry}
\usepackage{amsmath,amssymb,mathtools}
\usepackage[scr=rsfs,cal=euler]{mathalfa}
\usepackage{xfrac}
\usepackage[unicode,hidelinks,bookmarks=false]{hyperref}
\newcommand{\ie}{i.e.\ }
\newcommand{\cf}{cf.\ }
\newcommand{\resp}{respectively\ }
\newcommand{\Subsection}{\subsection}
\providecommand{\nobreakdash}{\nobreak\hbox{-}}
\setlength{\emergencystretch}{2em}
\multlinegap0pt
\usepackage{bm}
\usepackage{bbm}
\usepackage{dsfont}
\usepackage{xspace}
\usepackage{cancel}
\usepackage{mathtools}
\usepackage{amsthm}
\makeatletter
\@ifundefined{theorem}{\newtheorem{theorem}{Theorem}}{}
\@ifundefined{claim}{\newtheorem{claim}[theorem]{Claim}}{}
\makeatother
\def\emptyset{\mbox{{\rm \O}}}
\title[Decompositions of graphs with degree constraints relative to]{Decompositions of graphs with degree constraints relative to prescribed subgraphs}
\author{Peichao Wei, Muhuo Liu, Yang Wu, Zoran Stani\' c}
\date{Source: arXiv:2509.23078v1 (CC BY 4.0)}
\begin{document}
\maketitle

\noindent\emph{Source and licence.} Decompositions of graphs with degree constraints relative to prescribed subgraphs. Version arXiv:2509.23078v1 (CC BY 4.0), distributed under the Creative Commons Attribution 4.0 International licence, as recorded in the arXiv licence metadata for this exact version and shown on the abstract page https://arxiv.org/abs/2509.23078v1. Full rights notice: licenses/arxiv-2509.23078-CC-BY-4.0.txt.\par\smallskip

\noindent\textit{Editorial scope.} Counted unit: the Claim whose source label is \texttt{28c} in the article (source0.tex lines 294--296, source label \texttt{28c}), delivered together with the article's own complete original proof of it (source0.tex lines 297--321, one \texttt{proof} environment of 25 lines). No mathematics is written, repaired or re-derived: statement and proof are the article's own text line for line. Required context retained uncounted: 13t (lines 167--172); 23a (lines 218--220); 21c (lines 242--244); 222c (lines 256--258); 22c (lines 264--266). Every definition, display and label of those lines is retained unchanged. Omitted from this package and not redistributed: 1--166, 173--217, 221--241, 245--255, 259--263, 267--293, 322--510. Nothing inside a retained excerpt is omitted or altered: every retained body is delivered exactly as the article's lines are written, so no omission receipt of any kind accompanies this package. Citations: the retained excerpts carry no citation command at all, so no bibliography entry of the article is transcribed and no reference is redistributed. Reference adaptation: none was needed and none was made; every reference inside the delivered unit resolves inside it, so no unresolved reference or citation is produced. Printed slips kept literal: no printed word, symbol, hypothesis or number of the counted statement or of its proof was altered. Layout and class: the article's own class is \texttt{elsarticle} with packages including amssymb, amsmath, amsthm, latexsym, amsfonts, mathtools, graphicx, color; the wrapper is the lane's pinned \texttt{amsart} editorial wrapper at 10pt on A4, which restates the theorem-like environments the retained text needs and adds the article's own macro definitions that the retained text needs (\texttt{\textbackslash emptyset}), with extra wrapper packages (bm, bbm, dsfont, xspace, cancel, mathtools). The delivered numbering is the unit's own. Assets: no retained span contains a figure environment, a graphics-inclusion command, a PDF-page inclusion, a table or a TikZ picture; no image file is copied into this package, the package declares no asset dependency and no image file is redistributed with it. Licence: version arXiv:2509.23078v1 of this article is distributed under the Creative Commons Attribution 4.0 International licence, as recorded in the arXiv licence metadata for this exact version and shown on the abstract page https://arxiv.org/abs/2509.23078v1. A full rights notice is delivered at licenses/arxiv-2509.23078-CC-BY-4.0.txt. Characters: every retained excerpt is pure ASCII, so the delivered engine, XeLaTeX with its default Latin Modern text font, reports no missing character.
\begin{theorem}\label{13t}
	For a graph $ G $  with at least five vertices and $ a,b\colon V(G)\longrightarrow \mathbb{Z}_{\ge 0}$, suppose that  $ d_G(u)\ge a(u)+b(u)+h(u)$ holds  for $u\in T_{h(u)}(G)$, with $ h(u)\in \left\lbrace 0,1\right\rbrace$. Then $G$ admits an  $(a,b)$-feasible partition whenever one of the following holds:
	\begin{itemize} \item[(i)] $\min \left\lbrace a(u),b(u)\right\rbrace\ge 1-h(u)$ for $u\in T_{h(u)}(G)$ and $H$ is  obtained   by connecting a single vertex to exactly two vertices of $ K_{4} -e $;
		\item[(ii)] $\min \left\lbrace a(u),b(u)\right\rbrace\ge 2-h(u)$ for $u\in T_{h(u)}(G)$  and  $H\cong K_{2,3}$.
		\end{itemize}
\end{theorem}	

\begin{align}\label{23a}
	\omega(X_1\setminus\{u\},X_2\cup\{u\})-\omega(X_1,X_2)&=d_{X_2}(u)-d_{X_1}(u)+a(u)-b(u)\nonumber\\&\ge (b(u)+h(u))-a(u)+a(u)-b(u)\\&=h(u),\nonumber
\end{align}

	\begin{claim}\label{21c}
		For every vertex $ u\in V(G) $, $ d_G(u)=a(u)+b(u)+h(u) $.
	\end{claim}

\begin{claim}\label{222c}
	For every vertex $ u\in V(G) $, $\min\{a(u),b(u)\}\ge1$.
\end{claim}

\begin{claim}\label{22c}
		For every vertex $ u\in V(G) $, $  V(G)\setminus \left\lbrace u \right\rbrace  $ is  $ 2 $-good.
		\end{claim}

 \begin{claim}\label{28c}
If $ (X_1,X_2)\in\mathcal{P} $, then 	$\min\{ |X_1|,\left|X_2 \right|\} \ge2 $.
\end{claim}

\begin{proof} Let   $ u\in B_2(X_2) $. By Claim \ref{21c}, we have   $$ \left| X_1 \right|\ge e(u,X_1)\ge d_G(u)-(b(u)+h(u)-1)\ge a(u)+1\ge2.$$ 

To complete the argument it suffices to establish the same inequality for $|X_2|$.  Assume that $|X_2|=1$, i.e. $X_2=\{u\}$.  We begin by deriving a collection of structural constraints. Using these constraints, we then construct one or more $(a,b)$-feasible partitions which contradict the assumption, thereby excluding the case $|X_2|=1$.



If $\min \left\lbrace a(v),b(v)\right\rbrace\ge 2-h(v)$ holds for every $v\in X_1$, then  $$d_{X_1}(v)\ge a(v)+b(v)+h(v)-e(u,v)\ge a(v)+2-h(v)+h(v)-e(u,v)\ge a(v)+1,$$ for every vertex $ v\in X_1 $. However, this contradicts the assumption $(X_1,X_2)\in\mathcal{P}$. Therefore, the assumption (ii) of Theorem~\ref{13t} is not satisfied, meaning the defining graph of the set $T_1$, denoted by $H$, is as in (i): obtained   by connecting a single vertex to exactly two vertices of $ K_{4} -e $. 

   Since $ X_1 $ is $ 1 $-meager, there is a vertex $x\in B_1(X_1)$. Since  $ d_{X_2}(x)\ge b(x)+h(x)\ge 1 $, by Claim \ref{222c},  we have $ux\in E(G)$. If $ (X_1\setminus\{x\}, X_2\cup\{x\}) $ is a degenerate partition, then the inequality \eqref{23a} implies $$ \omega(X_1\setminus\{x\}, X_2\cup\{x\})-\omega(X_1,X_2)\ge h(x)\ge 0,$$ which contradicts  the minimality of $| X_1|$. Thus, the partition is not degenerate, which implies that $X_2\cup\{x\}$ is $2$-good, and so $1=d_{\{u,x\}}(u)\ge b(u)+h(u)$ and
$1=d_{\{u,x\}}(x)\ge b(x)+h(x)$. Hence, $b(u)=b(x)=1$ and $h(u)=h(x)=0$, by Claim \ref{222c}. This, in particular, means that $x,u\in T_0$. In addition, for every  $x\in B_1(X_1)$, from  $ux\in E(G)$  and $d_{\{u,x\}}(x)=b(x)=1$, by employing  Claim \ref{22c} we obtain $d_{X_1}(x)=a(x)$.  


From $x\in T_0$ and (again) $ux\in E(G)$, we have $$\left| N_{X_1}(x)\cap B_1(X_1)\right|\le 2,$$
as  $ \left| N_{X_1}(x)\cap B_1(X_1)\right|\ge 3 $ would imply $ x\in T_1(G) $.

 If $ \left| N_{X_1}(x_0)\cap B_1(X_1)\right|=0$ for some $x_0\in B_1(X_1)$, then $(X_1\setminus \{x_0\}, X_2\cup \{x_0\})$ is an $(a,b)$-feasible partition, which is impossible. We select $y_0\in N_{X_1}(x_0)\cap B_1(X_1)$.  By the symmetry of $x_0$ and $y_0$, we have $1\le  \left| N_{X_1}(y_0)\cap B_1(X_1)\right|\le 2$.

We claim that  $|N_{X_1}(y_0)\cap B_1(X_1)|=|N_{X_1}(x_0)\cap B_1(X_1)|=1$. The other possibility for these cardinalities is $2$. Without loss of generality, suppose that   $ N_{X_1}(y_0)\cap B_1(X_1)=\{x_0,w_0\} $.
Since $|V(G)|\ge 5$, we have  $X_1\ne \{w_0,y_0,x_0\}$.
Let $x\in X_1\setminus \{w_0,y_0,x_0\}$. If $x\in B_1(X_1)$, then $N_G(x)\cap \{w_0,y_0,x_0\}=\emptyset$ holds by the structure of $H$, as $x_0\in T_0$. If   $x\not\in B_1(X_1)$, then $|N_G(x)\cap \{w_0,y_0,x_0\}|\le 1$, by the same argument  and $d_{X_1}(x)\ge a(x)+1$. Let $X''_1=X_1\setminus \{w_0,y_0,v_0\}$ and $X''_2=X_2\cup \{w_0,y_0,x_0\}$. Irrespective of whether $x\in B_1(X_1)$, we have $d_{X''_1}(x)\ge a(x)$ and $d_{X''_2}(y)\ge 2>b(y)$ for every $y\in X''_2$. Thus, $(X''_1, X''_2)$ is an $(a, b) $-feasible partition, which is impossible. This confirms our claim.

  On the basis of $ |V(G)|\ge 5 $, let $x\in X_1\setminus\{x_0,y_0\}\neq\emptyset $.
  If $ \left| N(x)\cap \{x_0,y_0\}\right|\le 1$ holds  for every $x\in X_1\setminus\{x_0,y_0\}$, then     $(X_1\setminus\{x_0,y_0\}, X_2\cup\{x_0,y_0\})$ is an $ (a,b) $-feasible partition. Thus, there exists   $x\in X_1\setminus B_1(X_1)$ such that  $ \left| N(x)\cap \{x_0,y_0\}\right|=2$.   Since $u\in T_0$,  $x$ features as the unique vertex of  $ X_1\setminus B_1(X_1) $ with  $ \left| N(x)\cap \{x_0,y_0\}\right|=2$.

  Finally, if $ d_{X_1}(x)\ge a(x)+2 $, then $ (X_1\setminus\{x_0,y_0\}, X_2\cup\{x_0,y_0\}) $ is an $ (a,b) $-feasible partition. Hence, $d_{X_1}(x)=a(x)+1 $. From $ X_1\setminus\{x,y_0,x_0\}\neq\emptyset $, we have $ \left|N(v)\cap\{x,y_0,x_0\} \right|\le 1 $ for every $ v\in X_1\setminus\{x,y_0,x_0\}$, as $u\in T_0$. Hence, $ X_1\setminus\{x,y_0,x_0\} $ is $ a $-nice. It  is easy check  that $(X_1\setminus\{x,y_0,x_0\}, X_2\cup \{x,y_0,x_0\})$ features as an $(a,b)$-feasible partition, which is the final contradiction. \end{proof}

\end{document}
